
This work intersects four areas: quality-aware scaling laws, perplexity-based data selection, benchmark evaluation methodology, and model documentation standards.

\subsubsection{2.1 Data Quality in Pretraining and Scaling Laws}

The Chinchilla scaling law \citep{Hoffmann2022Training} characterizes loss as a function of model size N and training data volume D, assuming fixed data quality. Subramanyam et al. (2025) extend this framework to include a data quality scalar Q, yielding L(N, D, Q) = A/$N^alpha$ + B/($D^beta * Q^gamma$) + E with gamma\_CLM approximately 0.39. This formalization is important for understanding the role of data quality but operationally limited: Q must be measured through controlled experiments, not observationally from deployed models.

Myntti et al. (2025) show that text register (opinion, news, web) differentially affects benchmark scores. Diao et al. (2025) demonstrate that automated domain mixture optimization yields improvement on target-domain benchmarks. These studies establish that data composition affects benchmark performance but require corpus access to compute composition.

The present work operationalizes data quality through observable documentation indicators (binary presence/absence in model cards) rather than requiring corpus access. This approach scales to thousands of models but measures documentation rather than implementation -- a distinction discussed further in Section 6.

\subsubsection{2.2 Perplexity and Corpus-Level Quality Proxies}

Thrush et al. (2025) estimate perplexity-benchmark rank correlations across 90 Open LLM Leaderboard v1 models to identify high-quality data sources for pretraining. Shum et al. (2025) build on this with FastText-based data selection using pre-computed loss-benchmark correlations. Messmer et al. (2025) show that model-based filtering on MMLU can match rule-based filtering with 15\% of training tokens.

These methods require either inference access or actual training runs. The documentation proxy used here requires only model card scraping and scales to 4,493 models with no inference compute. This work uses Open LLM Leaderboard v2, which evaluates harder reasoning and mathematical capabilities (IFEval, BBH, MATH Lvl 5, GPQA, MUSR, MMLU-PRO) that differ from the v1 benchmark suite (ARC, HellaSwag, TruthfulQA, Winogrande, GSM8K, MMLU). Direct score comparisons with Thrush et al. (2025) are therefore not appropriate.

\subsubsection{2.3 Benchmark Evaluation and Contamination}

Wu et al. (2024) propose constructing contamination-free benchmarks from new knowledge. Lee et al. (2022) demonstrate that deduplication of pretraining data measurably improves downstream benchmark performance. These studies establish that documented decontamination and deduplication practices -- tracked here as binary features -- have real effects on benchmark outcomes.

\subsubsection{2.4 Model Documentation Standards}

Mitchell et al. (2019) introduce model cards as AI transparency infrastructure. Gebru et al. (2021) propose datasheets for datasets. Despite wide adoption on HuggingFace, model cards have primarily been studied qualitatively rather than for their relationship to independently measured outcomes. The present work is, to the authors' knowledge, the first to quantitatively link model card curation documentation to independently measured benchmark performance at the scale of thousands of models.

\subsubsection{2.5 Summary of Positioning}

\begin{table}[htbp]
\centering
\resizebox{\columnwidth}{!}{%
\begin{tabular}{llll}
\toprule
Work & Models & Documentation Signal & Performance Signal \\
\midrule
Thrush et al. (2025) & 90 & Computed (inference) & Open LLM Leaderboard v1 \\
Subramanyam et al. (2025) & Synthetic & Controlled (training) & Held-out loss \\
Myntti et al. (2025) & Few & Direct corpus access & Benchmark scores \\
This work & 4,493 & Human-written (cards) & Open LLM Leaderboard v2 \\
\bottomrule
\end{tabular}
}
\end{table}
